% =========================================================================
% Abstract + Introduction for the CD-GNN RA-L submission.
% Usage: abstract goes in \begin{abstract}...\end{abstract} in main.tex;
%        \input this file for the Introduction section.
% Numbers in [brackets] are provisional (from the 25-event 5-AMR runs of
% July 2026) -- regenerate from the final benchmark before submission.
% CLAIM GUARD: the "outperforms hand-crafted rules" claim requires the GNN
% retrained against the milk_run baseline to beat ~450 at m=5; until then,
% use the weaker phrasing marked ALT below.
% =========================================================================

% ---------------------------- ABSTRACT ----------------------------------
% Keep <= 200 words. Structure: context (1 sentence), gap (1-2), method
% (2-3), results (2), significance (1).
\newcommand{\CDGNNabstract}{%
Cross-docking terminals increasingly rely on fleets of autonomous mobile
robots (AMRs) to transfer unit loads from inbound docks to outbound
processing stations. Scheduling these transfers is a flexible job-shop
problem tightly coupled to the physical floor: docks are single-server
resources, robots consolidate loads under capacity limits, and every
trip must be routed collision-free through shared space. Existing
learned dispatchers idealize transport as fixed travel times, while
multi-agent path-finding methods route rigorously but ignore scheduling
structure. We present a graph-neural scheduling policy that closes this
gap: a sparse dynamic disjunctive graph encodes committed precedences,
congestion enters through calibrated features fitted against a
collision-aware executor, and the policy is trained with a stepwise
counterfactual baseline in which every decision is credited by
completing the schedule with a dispatching rule and executing it under
full collision-aware routing. On $60$-job instances with five AMRs, the
learned policy outperforms a tuned genetic algorithm by
[$6\%$] while constructing schedules [two orders of magnitude] faster,
and discovers load-consolidation behavior that its own training
baseline cannot express. A fleet-density study shows congestion costs
growing from [$2.6\%$] to [$12\%$] of makespan as the fleet doubles,
delimiting the regime where execution-aware scheduling pays.%
}

% ------------------------- I. INTRODUCTION ------------------------------

\section{Introduction}
\label{sec:intro}

Cross-docking replaces warehouse storage with synchronized flow:
freight arriving at inbound docks is transferred across the terminal
and dispatched from outbound doors within hours. As terminals adopt
fleets of autonomous mobile robots (AMRs) for this internal transfer,
the operator faces an integrated decision problem---which robot moves
which load, in what order---whose quality is measured not on a Gantt
chart but on a physical floor, where robots queue at single-server
docks, consolidate loads under capacity limits, and must never occupy
the same cell at the same time.

This problem is hard along two coupled axes. As a scheduling problem
it generalizes the flexible job shop with transportation: jobs are
pickup--delivery operation pairs over exclusive dock resources, and
vehicle assignment interacts with sequencing through carrying capacity.
As a robotics problem, realized travel and waiting times emerge from
collision-free multi-agent routing, so the makespan of a schedule is
defined only through its execution. Optimizing through this coupling
is what separates deployable schedulers from paper ones: we show
empirically that policies trained oblivious to execution density
degrade sharply when fleets densify
(Section~\ref{sec:experiments}).

Existing work addresses the two axes separately.
Learning-to-dispatch methods construct schedules with graph neural
policies \cite{zhang2020l2d,song2023fjsp} and have been extended to
transportation-aware variants \cite{moon2024hgs}, but transport is
reduced to fixed travel-time matrices---collisions, congestion, and
queueing are absent. Conversely, multi-agent path finding and its
pickup--delivery variants \cite{stern2019mapf,ma2017lifelong} route
rigorously through shared space but carry a deliberately thin task
model with neither dock exclusivity nor consolidation. A third line
integrates scheduling with conflict-free routing in container
terminals via hierarchical decomposition and per-instance search
\cite{zhong2020integrated,hierarchical2024agv}, without learning a
reusable policy. The intersection---a learned constructive scheduler
whose training signal reflects collision-aware execution---remains,
to our knowledge, unexplored.

We close this gap with \textsc{CD-GNN}, a graph-neural dispatcher for
AMR-served cross-docking. Three design decisions define it. First, the
policy plans on a lightweight surrogate whose travel and dock-wait
estimates are \emph{calibrated} offline against a collision-aware
executor, so the features every decision conditions on track physical
reality. Second, credit assignment is \emph{executor-grounded}: the
advantage of each decision is computed by completing the partial
schedule with a strong dispatching rule and executing both variants
under full collision-aware routing---per-step counterfactual credit
that no lower-bound reward provides. Third, the job graph is kept
\emph{sparse and dynamic}, encoding only committed precedences; we
show that the seemingly natural alternative---dense edges between jobs
sharing a dock---collapses job embeddings by over-smoothing and
prevents policy-gradient training from ever lifting off.

Our contributions are:
\begin{itemize}
\item a formulation of integrated dock-exclusive, capacity-coupled
  AMR scheduling as simulation-based optimization against a
  collision-aware executor, cast as a masked sequential decision
  process (Section~\ref{sec:problem});
\item \textsc{CD-GNN}, a $2.8\times10^{5}$-parameter,
  size- and fleet-agnostic dispatching policy, together with an
  offline calibration procedure that halves the surrogate--executor
  makespan gap (Sections~\ref{sec:architecture}--\ref{sec:calibration});
\item an executor-grounded stepwise counterfactual training scheme
  whose learned policies discover load consolidation---batching up to
  capacity---that the training baseline itself cannot express
  (Sections~\ref{sec:training}, \ref{sec:experiments});
\item an empirical study on $20$--$100$-job instances: the policy
  % ALT (pre-retrain phrasing): matches a tuned genetic algorithm at
  % two orders of magnitude less computation
  outperforms a tuned genetic algorithm by [$6\%$] at [two orders of
  magnitude] less computation, generalizes zero-shot across instance
  sizes, and a fleet-density analysis with a time-budget decomposition
  identifies when execution-aware training is necessary---and when it
  is not (Section~\ref{sec:experiments}).
\end{itemize}
